% ============================================================================
% CHAPTER 0: MATHEMATICAL FOUNDATIONS FOR CONTROL SYSTEMS
% ============================================================================

% ============================================================================
\chapter{Complex Numbers and Imaginary Arithmetic}
\label{ch:complex_numbers}
% ============================================================================

Complex numbers are fundamental to control systems analysis. They appear naturally in the solution of differential equations, in the Laplace transform, in frequency response analysis, and in stability analysis.

\section{Definition and Representation}

\begin{definitionbox}[title=Definition 0.1: Complex Number]
A \textbf{complex number} $z$ is defined as:
\begin{equation}
z = a + \jj b
\end{equation}
where $a, b \in \mathbb{R}$, and $\jj = \sqrt{-1}$ is the \textbf{imaginary unit}. The real part is $\Real(z) = a$ and the imaginary part is $\Imag(z) = b$.
\end{definitionbox}

\begin{note}
In electrical engineering, $\jj$ is often used instead of $\ii$ for the imaginary unit to avoid confusion with current. We adopt this convention throughout.
\end{note}

Complex numbers can be represented in multiple forms:

\subsubsection{Rectangular Form}
\begin{equation}
z = a + \jj b
\end{equation}

\subsubsection{Polar Form}
\begin{equation}
z = r(\cos\theta + \jj\sin\theta) = r \angle \theta
\end{equation}
where $r = |z| = \sqrt{a^2 + b^2}$ is the \textbf{magnitude} (or modulus) and $\theta = \arg(z) = \arctan(b/a)$ is the \textbf{argument} (or phase angle).

\subsubsection{Exponential Form (Euler's Formula)}
\begin{equation}
z = r \ee^{\jj\theta}
\end{equation}

Euler's formula, one of the most elegant results in mathematics, states:
\begin{equation}
\boxed{\ee^{\jj\theta} = \cos\theta + \jj\sin\theta}
\end{equation}

\begin{figure}[htbp]
\centering
\begin{tikzpicture}[scale=2]
    % Axes
    \draw[->] (-1.5,0) -- (2,0) node[right] {$\Real$};
    \draw[->] (0,-1.2) -- (0,1.5) node[above] {$\Imag$};

    % Grid
    \draw[help lines, gray!30] (-1.5,-1.2) grid (2,1.5);

    % Unit circle
    \draw[dashed,gray!60,thick] (0,0) circle (1cm);
    \node[below right,gray!60,font=\small] at (0.7,-0.7) {unit circle};

    % Complex number
    \coordinate (Z) at (1.3,0.9);
    \draw[->,thick,UofSCAtlantic] (0,0) -- (Z) node[above right,UofSCAtlantic,font=\bfseries] {$z = a + \jj b$};

    % Projections to axes (dashed red)
    \draw[dashed,UofSCGarnet!60,thick] (1.3,0) -- (Z);
    \draw[dashed,UofSCGarnet!60,thick] (0,0.9) -- (Z);

    % Projection points
    \draw[fill,UofSCGarnet] (1.3,0) circle (0.05cm);
    \draw[fill,UofSCGarnet] (0,0.9) circle (0.05cm);

    % Labels
    \node[below,UofSCGarnet,font=\bfseries] at (1.3,-0.15) {$a = \Real(z)$};
    \node[left,UofSCGarnet,font=\bfseries] at (-0.25,0.9) {$b = \Imag(z)$};

    % Angle (green)
    \draw[->,UofSCHorseshoe,thick] (0.5,0) arc (0:35:0.5) node[midway,right,UofSCHorseshoe,font=\small] {$\theta$};

    % Magnitude arc
    \draw[thick,UofSCAtlantic,opacity=0.5] (0,0) -- (0.65,0.45);
    \node[above left,UofSCAtlantic,font=\bfseries] at (0.65,0.45) {$r = |z|$};

    % Origin label
    \node[below left,font=\small] at (0,0) {$0$};
\end{tikzpicture}
\caption{Geometric representation of a complex number in the complex plane (Argand diagram). The blue vector shows the complex number $z$, with dashed red projections onto the real and imaginary axes. The green angle $\theta$ represents the argument, and the unit circle (dashed gray) provides a reference. Colors: Atlantic (blue) for magnitude, Garnet (red) for projections, Horseshoe (green) for angle.}
\label{fig:complex_plane}
\end{figure}

\begin{figure}[htbp]
\centering
\begin{tikzpicture}[scale=2]
    % Axes
    \draw[->] (-0.5,0) -- (2.5,0) node[right] {$\Real$};
    \draw[->] (0,-0.5) -- (0,2.2) node[above] {$\Imag$};

    % Grid
    \draw[help lines, gray!30] (-0.5,-0.5) grid (2.5,2.2);

    % Unit circle
    \draw[dashed,gray!60,thick] (0,0) circle (1cm);

    % Define coordinates
    % z1 = 0.8 * e^(j*25°) = 0.8(cos25 + j*sin25)
    \coordinate (Z1) at ({0.8*cos(25)},{0.8*sin(25)});
    % z2 = 1.2 * e^(j*40°) = 1.2(cos40 + j*sin40)
    \coordinate (Z2) at ({1.2*cos(40)},{1.2*sin(40)});
    % z1*z2 = 0.96 * e^(j*65°)
    \coordinate (Z1Z2) at ({0.96*cos(65)},{0.96*sin(65)});

    % Draw vectors
    \draw[->,thick,UofSCAtlantic] (0,0) -- (Z1) node[right,UofSCAtlantic,font=\small\bfseries] {$z_1 = r_1\ee^{\jj\theta_1}$};
    \draw[->,thick,UofSCHorseshoe] (0,0) -- (Z2) node[above right,UofSCHorseshoe,font=\small\bfseries] {$z_2 = r_2\ee^{\jj\theta_2}$};
    \draw[->,thick,UofSCGarnet,line width=1.5pt] (0,0) -- (Z1Z2) node[above,UofSCGarnet,font=\small\bfseries] {$z_1 z_2 = r_1 r_2\ee^{\jj(\theta_1+\theta_2)}$};

    % Angle arcs
    % theta1 = 25°
    \draw[->,UofSCAtlantic] (0.35,0) arc (0:25:0.35);
    \node[right,UofSCAtlantic,font=\scriptsize] at ({0.4*cos(12.5)},{0.4*sin(12.5)}) {$\theta_1$};

    % theta2 = 40° (from 0)
    \draw[->,UofSCHorseshoe] (0.5,0) arc (0:40:0.5);
    \node[right,UofSCHorseshoe,font=\scriptsize] at ({0.55*cos(20)},{0.55*sin(20)}) {$\theta_2$};

    % theta1 + theta2 = 65° (from 0)
    \draw[->,UofSCGarnet] (0.65,0) arc (0:65:0.65);
    \node[above right,UofSCGarnet,font=\scriptsize] at ({0.7*cos(32.5)},{0.7*sin(32.5)+0.05}) {$\theta_1+\theta_2$};

    % Magnitude labels with braces
    % r1 magnitude
    \draw[UofSCAtlantic,thick,decorate,decoration={brace,amplitude=3pt,mirror}]
        (0,-0.1) -- ({0.8*cos(25)},-0.1) node[midway,below=3pt,UofSCAtlantic,font=\scriptsize] {$r_1$};

    % r2 magnitude (along the vector direction, offset)
    \node[UofSCHorseshoe,font=\scriptsize,rotate=40] at ({0.6*cos(40)-0.1},{0.6*sin(40)+0.1}) {$r_2$};

    % r1*r2 magnitude
    \node[UofSCGarnet,font=\scriptsize,rotate=65] at ({0.48*cos(65)-0.12},{0.48*sin(65)+0.08}) {$r_1 r_2$};

    % Points at vector tips
    \draw[fill,UofSCAtlantic] (Z1) circle (0.04cm);
    \draw[fill,UofSCHorseshoe] (Z2) circle (0.04cm);
    \draw[fill,UofSCGarnet] (Z1Z2) circle (0.04cm);

    % Origin label
    \node[below left,font=\small] at (0,0) {$0$};

    % Key insight box
    \node[draw,thick,fill=white,font=\small,align=center] at (1.9,1.8) {
        Magnitudes multiply\\
        Angles add
    };
\end{tikzpicture}
\caption{Complex multiplication in polar form. When multiplying two complex numbers $z_1 = r_1\ee^{\jj\theta_1}$ and $z_2 = r_2\ee^{\jj\theta_2}$, the magnitudes multiply ($r_1 r_2$) and the angles add ($\theta_1 + \theta_2$). This geometric interpretation makes polar form ideal for multiplication and division operations.}
\label{fig:complex_multiplication}
\end{figure}

\begin{figure}[htbp]
\centering
\begin{tikzpicture}[scale=2]
    % Axes
    \draw[->] (-0.5,0) -- (2.5,0) node[right] {$\Real$};
    \draw[->] (0,-1.5) -- (0,1.5) node[above] {$\Imag$};

    % Grid
    \draw[help lines, gray!30] (-0.5,-1.5) grid (2.5,1.5);

    % Real axis highlight (reflection line)
    \draw[gray!60,thick,dashed] (-0.5,0) -- (2.5,0);
    \node[gray!60,font=\scriptsize,above] at (2.2,0.05) {reflection axis};

    % Define coordinates
    \coordinate (Z) at (1.5,0.9);
    \coordinate (Zconj) at (1.5,-0.9);
    \coordinate (Zmod2) at (3.06,0);  % |z|^2 = 1.5^2 + 0.9^2 = 3.06

    % Dashed reflection lines
    \draw[dashed,gray!50,thick] (Z) -- (Zconj);

    % Draw vectors for z and z*
    \draw[->,thick,UofSCAtlantic,line width=1.5pt] (0,0) -- (Z) node[above right,UofSCAtlantic,font=\bfseries] {$z = a + \jj b$};
    \draw[->,thick,UofSCHorseshoe,line width=1.5pt] (0,0) -- (Zconj) node[below right,UofSCHorseshoe,font=\bfseries] {$z^* = a - \jj b$};

    % Projections to axes for z
    \draw[dashed,UofSCGarnet!60,thick] (1.5,0) -- (Z);
    \draw[dashed,UofSCGarnet!60,thick] (0,0.9) -- (Z);

    % Projections to axes for z*
    \draw[dashed,UofSCGarnet!60,thick] (1.5,0) -- (Zconj);
    \draw[dashed,UofSCGarnet!60,thick] (0,-0.9) -- (Zconj);

    % Projection points on axes
    \draw[fill,UofSCGarnet] (1.5,0) circle (0.04cm);
    \draw[fill,UofSCGarnet] (0,0.9) circle (0.04cm);
    \draw[fill,UofSCGarnet] (0,-0.9) circle (0.04cm);

    % Labels for projections
    \node[below,UofSCGarnet,font=\small] at (1.5,-0.08) {$a$};
    \node[left,UofSCGarnet,font=\small] at (-0.05,0.9) {$b$};
    \node[left,UofSCGarnet,font=\small] at (-0.05,-0.9) {$-b$};

    % Angle arcs
    \draw[->,UofSCAtlantic] (0.4,0) arc (0:31:0.4);
    \node[right,UofSCAtlantic,font=\scriptsize] at (0.45,0.15) {$\theta$};

    \draw[->,UofSCHorseshoe] (0.4,0) arc (0:-31:0.4);
    \node[right,UofSCHorseshoe,font=\scriptsize] at (0.45,-0.15) {$-\theta$};

    % Points at vector tips
    \draw[fill,UofSCAtlantic] (Z) circle (0.04cm);
    \draw[fill,UofSCHorseshoe] (Zconj) circle (0.04cm);

    % Show |z|^2 on real axis (scaled to fit)
    \coordinate (Zmod2scaled) at (2.2,0);
    \draw[->,thick,UofSCGarnet,line width=1.5pt] (0,0) -- (Zmod2scaled);
    \draw[fill,UofSCGarnet] (Zmod2scaled) circle (0.04cm);
    \node[above,UofSCGarnet,font=\small\bfseries] at (2.2,0.08) {$z \cdot z^* = |z|^2$};

    % Magnitude labels
    \node[UofSCAtlantic,font=\scriptsize,rotate=31] at ({0.75*cos(31)+0.08},{0.75*sin(31)+0.12}) {$|z| = r$};
    \node[UofSCHorseshoe,font=\scriptsize,rotate=-31] at ({0.75*cos(31)+0.08},{-0.75*sin(31)-0.12}) {$|z^*| = r$};

    % Origin label
    \node[below left,font=\small] at (0,0) {$0$};

    % Key property box
    \node[draw,thick,fill=white,font=\small,align=center] at (0.6,1.25) {
        $z^* = $ reflection of $z$\\
        across real axis
    };
\end{tikzpicture}
\caption{Complex conjugate visualization. The conjugate $z^* = a - \jj b$ is the reflection of $z = a + \jj b$ across the real axis. Both have the same magnitude $|z| = |z^*| = r$, but opposite angles $\pm\theta$. The product $z \cdot z^* = |z|^2 = a^2 + b^2$ is always a real, non-negative number lying on the real axis.}
\label{fig:complex_conjugate}
\end{figure}

\begin{figure}[htbp]
\centering
\begin{tikzpicture}[scale=2.5]
    % Axes
    \draw[->] (-1.4,0) -- (1.4,0) node[right] {$\Real$};
    \draw[->] (0,-1.4) -- (0,1.4) node[above] {$\Imag$};

    % Grid
    \draw[help lines, gray!30] (-1.4,-1.4) grid (1.4,1.4);

    % Unit circle
    \draw[thick,gray!60] (0,0) circle (1cm);
    \node[below right,gray!60,font=\scriptsize] at (0.7,-0.75) {unit circle};

    % Define the 5th roots of unity: e^(j*2*pi*k/5) for k=0,1,2,3,4
    % Angles: 0, 72, 144, 216, 288 degrees
    \coordinate (W0) at ({cos(0)},{sin(0)});           % k=0: 1
    \coordinate (W1) at ({cos(72)},{sin(72)});         % k=1
    \coordinate (W2) at ({cos(144)},{sin(144)});       % k=2
    \coordinate (W3) at ({cos(216)},{sin(216)});       % k=3
    \coordinate (W4) at ({cos(288)},{sin(288)});       % k=4

    % Draw pentagon (connecting lines)
    \draw[thick,UofSCGarnet!50] (W0) -- (W1) -- (W2) -- (W3) -- (W4) -- cycle;

    % Draw vectors from origin to each root
    \draw[->,thick,UofSCAtlantic] (0,0) -- (W0);
    \draw[->,thick,UofSCAtlantic] (0,0) -- (W1);
    \draw[->,thick,UofSCAtlantic] (0,0) -- (W2);
    \draw[->,thick,UofSCAtlantic] (0,0) -- (W3);
    \draw[->,thick,UofSCAtlantic] (0,0) -- (W4);

    % Points at each root
    \draw[fill,UofSCGarnet] (W0) circle (0.05cm);
    \draw[fill,UofSCGarnet] (W1) circle (0.05cm);
    \draw[fill,UofSCGarnet] (W2) circle (0.05cm);
    \draw[fill,UofSCGarnet] (W3) circle (0.05cm);
    \draw[fill,UofSCGarnet] (W4) circle (0.05cm);

    % Labels for each root
    \node[right,UofSCGarnet,font=\small\bfseries] at ({1.08*cos(0)},{1.08*sin(0)}) {$\omega_0 = 1$};
    \node[above right,UofSCGarnet,font=\small\bfseries] at ({1.08*cos(72)},{1.08*sin(72)}) {$\omega_1$};
    \node[above left,UofSCGarnet,font=\small\bfseries] at ({1.08*cos(144)},{1.08*sin(144)}) {$\omega_2$};
    \node[below left,UofSCGarnet,font=\small\bfseries] at ({1.08*cos(216)},{1.08*sin(216)}) {$\omega_3$};
    \node[below right,UofSCGarnet,font=\small\bfseries] at ({1.08*cos(288)},{1.08*sin(288)}) {$\omega_4$};

    % k-value labels (smaller, offset)
    \node[right,font=\scriptsize,gray!70] at ({1.25*cos(0)},{1.25*sin(0)-0.08}) {$k=0$};
    \node[above right,font=\scriptsize,gray!70] at ({1.25*cos(72)+0.02},{1.25*sin(72)}) {$k=1$};
    \node[above left,font=\scriptsize,gray!70] at ({1.25*cos(144)-0.02},{1.25*sin(144)}) {$k=2$};
    \node[below left,font=\scriptsize,gray!70] at ({1.25*cos(216)-0.02},{1.25*sin(216)}) {$k=3$};
    \node[below right,font=\scriptsize,gray!70] at ({1.25*cos(288)+0.02},{1.25*sin(288)}) {$k=4$};

    % Angle arcs for first two roots
    \draw[->,UofSCHorseshoe,thick] (0.25,0) arc (0:72:0.25);
    \node[above right,UofSCHorseshoe,font=\scriptsize] at (0.2,0.15) {$\frac{2\pi}{5}$};

    % Origin label
    \node[below left,font=\small] at (0,0) {$0$};

    % Axis tick marks
    \draw[thick] (1,-0.05) -- (1,0.05) node[below=3pt,font=\scriptsize] {$1$};
    \draw[thick] (-1,-0.05) -- (-1,0.05) node[below=3pt,font=\scriptsize] {$-1$};
    \draw[thick] (-0.05,1) -- (0.05,1) node[left=3pt,font=\scriptsize] {$\jj$};
    \draw[thick] (-0.05,-1) -- (0.05,-1) node[left=3pt,font=\scriptsize] {$-\jj$};

    % Formula box
    \node[draw,thick,fill=white,font=\small,align=center] at (0,-1.25) {
        $\omega_k = \ee^{\jj 2\pi k/5}$, \quad $k = 0, 1, 2, 3, 4$
    };
\end{tikzpicture}
\caption{The fifth roots of unity. The $n$th roots of unity are $n$ complex numbers evenly distributed on the unit circle, satisfying $z^n = 1$. For $n=5$, these roots $\omega_k = \ee^{\jj 2\pi k/5}$ are separated by angles of $2\pi/5 = 72°$ and form a regular pentagon. The root $\omega_0 = 1$ always lies on the positive real axis.}
\label{fig:roots_of_unity}
\end{figure}


\section{Arithmetic Operations}

\subsubsection{Addition and Subtraction}
\begin{equation}
(a + \jj b) \pm (c + \jj d) = (a \pm c) + \jj(b \pm d)
\end{equation}

\subsubsection{Multiplication}
\begin{equation}
(a + \jj b)(c + \jj d) = (ac - bd) + \jj(ad + bc)
\end{equation}
In polar form: $r_1 \ee^{\jj\theta_1} \cdot r_2 \ee^{\jj\theta_2} = r_1 r_2 \ee^{\jj(\theta_1 + \theta_2)}$

\subsubsection{Division}
\begin{equation}
\frac{a + \jj b}{c + \jj d} = \frac{(a + \jj b)(c - \jj d)}{(c + \jj d)(c - \jj d)} = \frac{(ac + bd) + \jj(bc - ad)}{c^2 + d^2}
\end{equation}
In polar form: $\frac{r_1 \ee^{\jj\theta_1}}{r_2 \ee^{\jj\theta_2}} = \frac{r_1}{r_2} \ee^{\jj(\theta_1 - \theta_2)}$

\subsubsection{Complex Conjugate}
\begin{equation}
z^* = \overline{z} = a - \jj b
\end{equation}

Important property: $z \cdot z^* = |z|^2 = a^2 + b^2$

\subsubsection{Powers (De Moivre's Theorem)}
\begin{equation}
z^n = r^n \ee^{\jj n\theta} = r^n(\cos n\theta + \jj\sin n\theta)
\end{equation}

\subsubsection{Roots}
The $n$th roots of $z = r\ee^{\jj\theta}$ are:
\begin{equation}
z^{1/n} = r^{1/n} \ee^{\jj(\theta + 2\pi k)/n}, \quad k = 0, 1, \ldots, n-1
\end{equation}

\begin{examplebox}[title=Example 0.1: Complex Number Operations]
Find the polar form, magnitude, and phase of $z = 3 + 4\jj$.

\textbf{Solution:}
\begin{align}
|z| &= \sqrt{3^2 + 4^2} = \sqrt{25} = 5 \\
\theta &= \arctan\left(\frac{4}{3}\right) \approx 53.13^{\circ} = 0.927 \text{ rad} \\
z &= 5\ee^{\jj(0.927)} = 5\angle 53.13^{\circ}
\end{align}
\end{examplebox}

\begin{examplebox}[title=Example 0.2: Multiplication and Division in Polar Form]
Given $z_1 = 2\ee^{\jj\pi/6}$ and $z_2 = 3\ee^{\jj\pi/4}$, compute $z_1 \cdot z_2$ and $z_1 / z_2$.

\textbf{Solution:}

For multiplication, we multiply magnitudes and add angles:
\begin{align}
z_1 \cdot z_2 &= (2)(3) \ee^{\jj(\pi/6 + \pi/4)} = 6\ee^{\jj(5\pi/12)}
\end{align}
Converting to rectangular form: $z_1 \cdot z_2 = 6\cos(5\pi/12) + 6\jj\sin(5\pi/12) \approx 1.55 + 5.80\jj$.

For division, we divide magnitudes and subtract angles:
\begin{align}
\frac{z_1}{z_2} &= \frac{2}{3} \ee^{\jj(\pi/6 - \pi/4)} = \frac{2}{3}\ee^{-\jj\pi/12}
\end{align}
Converting to rectangular form: $z_1/z_2 = \frac{2}{3}\cos(-\pi/12) + \frac{2}{3}\jj\sin(-\pi/12) \approx 0.644 - 0.173\jj$.

\textbf{Key insight:} Polar form makes multiplication and division much simpler than rectangular form, which is why it is preferred in frequency response analysis.
\end{examplebox}

\begin{examplebox}[title=Example 0.3: Finding All Cube Roots of a Complex Number]
Find all cube roots of $z = 8$.

\textbf{Solution:}

First, write $z = 8$ in polar form: $z = 8\ee^{\jj \cdot 0} = 8\ee^{\jj \cdot 2\pi k}$ for any integer $k$.

The $n$th roots of $z = r\ee^{\jj\theta}$ are given by:
\begin{equation*}
z^{1/n} = r^{1/n} \ee^{\jj(\theta + 2\pi k)/n}, \quad k = 0, 1, \ldots, n-1
\end{equation*}

For the cube roots ($n = 3$) of $8\ee^{\jj \cdot 0}$:
\begin{align}
z_0 &= 8^{1/3} \ee^{\jj(0 + 2\pi \cdot 0)/3} = 2\ee^{\jj \cdot 0} = 2 \\
z_1 &= 8^{1/3} \ee^{\jj(0 + 2\pi \cdot 1)/3} = 2\ee^{\jj 2\pi/3} = 2\left(-\frac{1}{2} + \jj\frac{\sqrt{3}}{2}\right) = -1 + \jj\sqrt{3} \\
z_2 &= 8^{1/3} \ee^{\jj(0 + 2\pi \cdot 2)/3} = 2\ee^{\jj 4\pi/3} = 2\left(-\frac{1}{2} - \jj\frac{\sqrt{3}}{2}\right) = -1 - \jj\sqrt{3}
\end{align}

\textbf{Verification:} We can verify $(2)^3 = 8$, $(-1 + \jj\sqrt{3})^3 = 8$, and $(-1 - \jj\sqrt{3})^3 = 8$.

\textbf{Geometric interpretation:} The three roots are equally spaced on a circle of radius $2$ in the complex plane, separated by $120^{\circ}$ (i.e., $2\pi/3$ radians).
\end{examplebox}

\begin{examplebox}[title=Example 0.4: Deriving Trigonometric Identities Using Euler's Formula]
Use Euler's formula to derive the double-angle formulas for $\cos(2\theta)$ and $\sin(2\theta)$.

\textbf{Solution:}

From Euler's formula: $\ee^{\jj\theta} = \cos\theta + \jj\sin\theta$.

Consider $\ee^{\jj(2\theta)}$. We can write this in two ways:

\textbf{Method 1:} Direct application of Euler's formula:
\begin{equation*}
\ee^{\jj(2\theta)} = \cos(2\theta) + \jj\sin(2\theta)
\end{equation*}

\textbf{Method 2:} Using properties of exponents:
\begin{align*}
\ee^{\jj(2\theta)} &= \left(\ee^{\jj\theta}\right)^2 = (\cos\theta + \jj\sin\theta)^2 \\
&= \cos^2\theta + 2\jj\cos\theta\sin\theta + \jj^2\sin^2\theta \\
&= (\cos^2\theta - \sin^2\theta) + \jj(2\cos\theta\sin\theta)
\end{align*}

Equating real and imaginary parts from both methods:
\begin{align}
\Real: \quad \cos(2\theta) &= \cos^2\theta - \sin^2\theta \\
\Imag: \quad \sin(2\theta) &= 2\cos\theta\sin\theta
\end{align}

\textbf{Extension:} This technique can be used to derive any multiple-angle formula. For example, expanding $(\ee^{\jj\theta})^3$ yields the triple-angle formulas.
\end{examplebox}

\section{Properties of Complex Operations}

This section provides formal proofs of fundamental properties of complex number operations. These properties are essential for understanding frequency response analysis and stability theory.

\begin{keyconceptbox}[title=Key Property: Magnitude of a Product]
For any two complex numbers $z_1$ and $z_2$:
\begin{equation}
|z_1 \cdot z_2| = |z_1| \cdot |z_2|
\end{equation}
The magnitude of a product equals the product of the magnitudes.
\end{keyconceptbox}

\begin{proof}
Let $z_1 = r_1 \ee^{\jj\theta_1}$ and $z_2 = r_2 \ee^{\jj\theta_2}$ be two complex numbers in polar form, where $r_1 = |z_1|$ and $r_2 = |z_2|$.

The product is:
\begin{equation*}
z_1 \cdot z_2 = r_1 \ee^{\jj\theta_1} \cdot r_2 \ee^{\jj\theta_2} = r_1 r_2 \ee^{\jj(\theta_1 + \theta_2)}
\end{equation*}

Taking the magnitude:
\begin{equation*}
|z_1 \cdot z_2| = |r_1 r_2 \ee^{\jj(\theta_1 + \theta_2)}| = r_1 r_2 = |z_1| \cdot |z_2|
\end{equation*}

where we used the fact that $|\ee^{\jj\phi}| = 1$ for any real $\phi$.
\end{proof}

\begin{keyconceptbox}[title=Key Property: Argument of a Product]
For any two complex numbers $z_1$ and $z_2$ (with $z_1, z_2 \neq 0$):
\begin{equation}
\arg(z_1 \cdot z_2) = \arg(z_1) + \arg(z_2) \pmod{2\pi}
\end{equation}
The argument of a product equals the sum of the arguments (modulo $2\pi$).
\end{keyconceptbox}

\begin{proof}
Let $z_1 = r_1 \ee^{\jj\theta_1}$ and $z_2 = r_2 \ee^{\jj\theta_2}$, where $\theta_1 = \arg(z_1)$ and $\theta_2 = \arg(z_2)$.

The product is:
\begin{equation*}
z_1 \cdot z_2 = r_1 r_2 \ee^{\jj(\theta_1 + \theta_2)}
\end{equation*}

By definition of the argument (the angle in polar form):
\begin{equation*}
\arg(z_1 \cdot z_2) = \theta_1 + \theta_2 = \arg(z_1) + \arg(z_2) \pmod{2\pi}
\end{equation*}

The modulo $2\pi$ accounts for the fact that angles differing by $2\pi$ represent the same direction.
\end{proof}

\begin{keyconceptbox}[title=Key Property: Magnitude Squared via Conjugate]
For any complex number $z$:
\begin{equation}
|z|^2 = z \cdot z^*
\end{equation}
The squared magnitude equals the product of a number and its complex conjugate.
\end{keyconceptbox}

\begin{proof}
Let $z = a + \jj b$ where $a, b \in \mathbb{R}$.

The complex conjugate is $z^* = a - \jj b$.

Computing the product:
\begin{align*}
z \cdot z^* &= (a + \jj b)(a - \jj b) \\
&= a^2 - \jj ab + \jj ab - \jj^2 b^2 \\
&= a^2 - (-1)b^2 \\
&= a^2 + b^2
\end{align*}

By definition of magnitude: $|z|^2 = (\sqrt{a^2 + b^2})^2 = a^2 + b^2$.

Therefore: $|z|^2 = z \cdot z^*$.

\textbf{Alternative proof using polar form:} Let $z = r\ee^{\jj\theta}$. Then $z^* = r\ee^{-\jj\theta}$, and:
\begin{equation*}
z \cdot z^* = r\ee^{\jj\theta} \cdot r\ee^{-\jj\theta} = r^2 \ee^{\jj \cdot 0} = r^2 = |z|^2
\end{equation*}
\end{proof}

\begin{note}
The property $|z|^2 = z \cdot z^*$ is particularly useful in control systems when computing power spectral densities and when rationalizing complex-valued transfer functions for frequency response analysis.
\end{note}

\section{Complex Numbers in Control Systems}

Complex numbers appear throughout control theory:

\begin{itemize}
    \item \textbf{Poles and Zeros:} The roots of transfer function polynomials are generally complex.
    \item \textbf{Eigenvalues:} System matrices often have complex eigenvalues that determine stability and oscillatory behavior.
    \item \textbf{Frequency Response:} Evaluating transfer functions at $s = \jj\omega$ yields complex numbers representing magnitude and phase.
    \item \textbf{Sinusoidal Steady-State:} Complex exponentials simplify analysis of sinusoidal inputs.
\end{itemize}

% ============================================================================
